\section{Propagation of Alternative Equations to Locked Diagnostics}
\label{app:alternative-propagation}

The following out-of-class controls verify that empirical-residual ranking is not used in the proof.  We propagated representative comparison values through the same equations and applied the locked PR-III references only after generation.  ``Correct'' in this appendix means inside the diagnostic interval accepted by PR-III; it does not mean exact equality.

\begin{table}[H]
\centering
\small
\caption{Compact-branch propagation to the locked rubidium
$\alpha^{-1}=137.035999206(11)$ diagnostic.}
\label{tab:compact-diagnostic-propagation}
\begin{tabular}{@{}lrr@{}}
\toprule
Branch & Generated $\alpha^{-1}$ & Distance \\
\midrule
Released preprint $(-1,\Pi_{13})$ & 137.035999207530239 & $0.1391\sigma$ \\
Alternative $(0,\Pi_{13})$ & 137.035999207470098 & $0.1336\sigma$ \\
Alternative $(+1,\Pi_{13})$ & 137.035999207409986 & $0.1282\sigma$ \\
Comparison $(-1,\Pi_{15})$ & 254.208754866216196 & outside interval \\
Comparison $(-1,\Pi_{17})$ & 384.885897406270487 & outside interval \\
\bottomrule
\end{tabular}
\end{table}

The three $\Pi_{13}$ branches are exactly distinct, but all lie inside the same one-standard-deviation interval.  They are statistically indistinguishable at this gate, so residual ordering is deliberately not used.  The exact signed boundary-incidence theorem supplies the selection.

For the PR-III electromagnetic correction, none of the 21 alternative bounded-degree monomials passes the rubidium one-sigma gate when used as a single unit-coefficient replacement term.  The PR-III $D_1+D_2$ sum does pass.  This numerical result is corroborative only.  The exponent rule is a fixed inherited PR-III finite-tier composition rule, and the 21 replacement monomials lie outside that declared grammar.

The electroweak coefficient propagation uses
\begin{equation}
 \Pi_W=\frac{a}{N_{\rm gen}}\epsilon_V,\qquad
 \Pi_Z=\frac{b}{N_{\rm gen}}\epsilon_V+
       \frac{k}{N_{\rm gen}}\epsilon_V^2.
\end{equation}
Besides the released-preprint $(a,b,k)=(2,-1,2)$, two alternatives pass both locked one-sigma mass gates:
\begin{table}[H]
\centering
\small
\caption{Electroweak coefficient triplets passing both locked mass gates.}
\label{tab:ew-alternative-diagnostics}
\begin{tabular}{@{}lrrrr@{}}
\toprule
$(a,b,k)$ & $M_W$ (GeV) & $M_Z$ (GeV) & $M_W$ distance & $M_Z$ distance \\
\midrule
$(2,-1,1)$ & 80.37394285 & 91.18590923 & $0.3566\sigma$ & $0.8051\sigma$ \\
$(2,-1,2)$ released preprint & 80.37394285 & 91.18767643 & $0.3566\sigma$ & $0.0364\sigma$ \\
$(2,-1,3)$ & 80.37394285 & 91.18944360 & $0.3566\sigma$ & $0.8779\sigma$ \\
\bottomrule
\end{tabular}
\end{table}

All three coefficients lie inside the stated diagnostic gate, so that gate is corroborative rather than generative.  The canonical normalized-response theorem supplies the exact selection $k=2$ independently of residual ranking.
